% GENERATED FILE. Edit canonical lessons, then run npm run generate:books.
% canonical-source-hash: fnv1a64:7fecf419b08a8488
% canonical-lessons: ES-C677-monumento, ES-C677-catedral, ES-C677-castillo, ES-C677-palacio, ES-C677-muralla

\chapter{Monumentos}
\label{ch:es-monumentos}
\hlchaptermodality{677}

\begin{chapteropening}
\textbf{By the end of this chapter:} \emph{I can talk about the sights of an old town: monuments, the cathedral, the castle, the palace and the walls.}

Complete the last lesson of chapter 677: I can talk about the sights of an old town: monuments, the cathedral, the castle, the palace and the walls.
\end{chapteropening}

\section[el monumento]{el monumento — the monument, the sight}

\label{lesson:ES-C677-monumento}



\begin{quote}
Before the new one: say the Spanish for the leaflet, then the Spanish for to haggle.
\end{quote}

\subsection*{You'll want to know: el monumento}
\textbf{el monumento} — \textquotedblleft{}the monument, the sight\textquotedblright{}.

\begin{quote}
Visitamos el monumento más famoso.
\end{quote}

\begin{quote}
\emph{We visited the most famous monument.}
\end{quote}

\emph{Monumentos} are the sights of a town.

\begin{cousinweb}
\emph{monumento} is a learned borrowing from Latin \emph{monumentum} \textquotedblleft{}that which keeps a memory, a memorial\textquotedblright{} (Lewis \& Short; EtymDB), from \emph{monēre} \textquotedblleft{}to remind\textquotedblright{}.
\end{cousinweb}

\begin{grammarlens}[title={What you've built}]
One more word for telling, explaining and describing at B1.
\end{grammarlens}

\subsection*{Your turn}
\begin{itemize}
\raggedright
  \item \emph{Say it:} \emph{el monumento}
  \item \emph{Say it:} \emph{el monumento}, once more
  \item \emph{Say it:} a sentence of your own with \emph{el monumento}
  \item \emph{Recall:} say the Spanish for the leaflet, then the Spanish for to haggle, then say \emph{el monumento} again
\end{itemize}

\begin{tcolorbox}[breakable,colback=teal!4,colframe=teal!35!black,title={Before you move on}]
What does el monumento mean? (\textquotedblleft{}The monument, the sight\textquotedblright{}.) Say it once more.
\end{tcolorbox}
\section[la catedral]{la catedral — the cathedral}

\label{lesson:ES-C677-catedral}



\begin{quote}
Before the new one: say the Spanish for to haggle, then the Spanish for the monument.
\end{quote}

\subsection*{You'll want to know: la catedral}
\textbf{la catedral} — \textquotedblleft{}the cathedral\textquotedblright{}.

\begin{quote}
La catedral es del siglo XIII.
\end{quote}

\begin{quote}
\emph{The cathedral is from the thirteenth century.}
\end{quote}

The bishop’s church of a city.

\begin{cousinweb}
\emph{catedral} is a learned borrowing from Latin \emph{(ecclesia) cathedrālis} \textquotedblleft{}the church with the chair\textquotedblright{}, from \emph{cathedra} \textquotedblleft{}a chair\textquotedblright{} (Lewis \& Short), Greek \emph{kathédra}: the bishop’s seat.
\end{cousinweb}

\begin{grammarlens}[title={What you've built}]
One more word for telling, explaining and describing at B1.
\end{grammarlens}

\subsection*{Your turn}
\begin{itemize}
\raggedright
  \item \emph{Say it:} \emph{la catedral}
  \item \emph{Say it:} \emph{la catedral}, once more
  \item \emph{Say it:} a sentence of your own with \emph{la catedral}
  \item \emph{Recall:} say the Spanish for to haggle, then the Spanish for the monument, then say \emph{la catedral} again
\end{itemize}

\begin{tcolorbox}[breakable,colback=teal!4,colframe=teal!35!black,title={Before you move on}]
What does la catedral mean? (\textquotedblleft{}The cathedral\textquotedblright{}.) Say it once more.
\end{tcolorbox}
\section[el castillo]{el castillo — the castle}

\label{lesson:ES-C677-castillo}



\begin{quote}
Before the new one: say the Spanish for the monument, then the Spanish for the cathedral.
\end{quote}

\subsection*{You'll want to know: el castillo}
\textbf{el castillo} — \textquotedblleft{}the castle\textquotedblright{}.

\begin{quote}
Desde el castillo se ve todo el pueblo.
\end{quote}

\begin{quote}
\emph{The whole town can be seen from the castle.}
\end{quote}

\emph{Hacer castillos en el aire} — to build castles in the air.

\begin{cousinweb}
\emph{castillo} is inherited from Latin \emph{castellum} \textquotedblleft{}a castle, a fort\textquotedblright{} (Lewis \& Short), a little \emph{castrum} \textquotedblleft{}fortified camp\textquotedblright{}.
\end{cousinweb}

\begin{grammarlens}[title={What you've built}]
One more word for telling, explaining and describing at B1.
\end{grammarlens}

\subsection*{Your turn}
\begin{itemize}
\raggedright
  \item \emph{Say it:} \emph{el castillo}
  \item \emph{Say it:} \emph{el castillo}, once more
  \item \emph{Say it:} a sentence of your own with \emph{el castillo}
  \item \emph{Recall:} say the Spanish for the monument, then the Spanish for the cathedral, then say \emph{el castillo} again
\end{itemize}

\begin{tcolorbox}[breakable,colback=teal!4,colframe=teal!35!black,title={Before you move on}]
What does el castillo mean? (\textquotedblleft{}The castle\textquotedblright{}.) Say it once more.
\end{tcolorbox}
\section[el palacio]{el palacio — the palace}

\label{lesson:ES-C677-palacio}



\begin{quote}
Before the new one: say the Spanish for the cathedral, then the Spanish for the castle.
\end{quote}

\subsection*{You'll want to know: el palacio}
\textbf{el palacio} — \textquotedblleft{}the palace\textquotedblright{}.

\begin{quote}
El rey vivía en este palacio.
\end{quote}

\begin{quote}
\emph{The king lived in this palace.}
\end{quote}

\emph{El Palacio de Justicia} — the law courts.

\begin{cousinweb}
\emph{palacio} is a learned borrowing from Latin \emph{palātium} (EtymDB), named after the Palatine hill in Rome, where the emperors had their houses.
\end{cousinweb}

\begin{grammarlens}[title={What you've built}]
One more word for telling, explaining and describing at B1.
\end{grammarlens}

\subsection*{Your turn}
\begin{itemize}
\raggedright
  \item \emph{Say it:} \emph{el palacio}
  \item \emph{Say it:} \emph{el palacio}, once more
  \item \emph{Say it:} a sentence of your own with \emph{el palacio}
  \item \emph{Recall:} say the Spanish for the cathedral, then the Spanish for the castle, then say \emph{el palacio} again
\end{itemize}

\begin{tcolorbox}[breakable,colback=teal!4,colframe=teal!35!black,title={Before you move on}]
What does el palacio mean? (\textquotedblleft{}The palace\textquotedblright{}.) Say it once more.
\end{tcolorbox}
\section[la muralla]{la muralla — the city wall, the ramparts}

\label{lesson:ES-C677-muralla}



\begin{quote}
Before the new one: say the Spanish for the castle, then the Spanish for the palace.
\end{quote}

\subsection*{You'll want to know: la muralla}
\textbf{la muralla} — \textquotedblleft{}the city wall, the ramparts\textquotedblright{}.

\begin{quote}
La muralla rodea el casco antiguo.
\end{quote}

\begin{quote}
\emph{The city wall runs round the old town.}
\end{quote}

A garden wall is \emph{un muro}; a city wall is \emph{la muralla}.

\begin{cousinweb}
\emph{muralla} is built on \emph{muro} \textquotedblleft{}wall\textquotedblright{}, which is inherited from Latin \emph{mūrus} (Lewis \& Short).
\end{cousinweb}

\begin{grammarlens}[title={What you've built}]
One more word for telling, explaining and describing at B1.
\end{grammarlens}

\subsection*{Your turn}
\begin{itemize}
\raggedright
  \item \emph{Say it:} \emph{la muralla}
  \item \emph{Say it:} \emph{la muralla}, once more
  \item \emph{Say it:} a sentence of your own with \emph{la muralla}
  \item \emph{Recall:} say the Spanish for the castle, then the Spanish for the palace, then say \emph{la muralla} again
\end{itemize}

\begin{tcolorbox}[breakable,colback=teal!4,colframe=teal!35!black,title={Before you move on}]
What does la muralla mean? (\textquotedblleft{}The city wall, the ramparts\textquotedblright{}.) Say it once more.
\end{tcolorbox}
